\section{Environment}
\label{sec:environment}

\subsection{Hardware Environment}

The experiments will run on Google Colab (free tier), using a single NVIDIA T4.

\subsection{Software Environment}

\begin{itemize}
    \item \textbf{Language:} Python 3.10+
    \item \textbf{Frameworks:} PyTorch, Hugging Face \texttt{transformers}
    \item \textbf{Inference Optimization:} \texttt{vLLM} for fast batch inference to minimize compute time.
\end{itemize}

\subsection{Dataset Requirements}

No model training is required. Instead, an evaluation dataset of 500 examples will be created from the multilingual MKQA or TyDi QA benchmarks. The dataset pairs English and Korean claims inserted into simple, consistent sentence templates to prevent structural differences. To ensure translations convey the same meaning, semantic consistency will be checked using back-translation and embedding cosine similarity.

\subsection{Model Selection}

\begin{itemize}
    \item \textbf{Qwen-2.5-7B-Instruct:} Chosen for its strong multilingual performance and a size suitable for a single T4 GPU.
    \item \textbf{Llama-3-8B-Instruct:} Chosen as a widely used baseline with well-documented English-centric bias.
\end{itemize}

\subsection{Experimental Pipeline}

\begin{enumerate}
    \item \textit{Data preparation:} Extract the question and correct answer, then generate a plausible false answer.
    \item \textit{Templating:} Insert answers into standardized English and Korean sentence templates.
    \item \textit{Prompting:} Present the model with the question alongside the two conflicting claims.
    \item \textit{Inference:} Generate answers using greedy decoding (temperature = 0) for reproducible outputs.
    \item \textit{Evaluation:} Determine which claim the model chose using an LLM-as-a-judge, avoiding the fragility of exact string matching.
\end{enumerate}